% Bewertungspaket zum Gesamttest «Bruchterme, Gleichungen und Formeln» — Quelle des PDFs (python3 scripts/build-lp-pdf.py gleichungen-formeln).
% Ganze Punkte je Teilschritt; Werte und Fehlerwerte mit python3 nachgerechnet (scripts/lp/gleichungen-formeln/zahlen.py).
% Gerüst aus den Bewertungspaketen Bruchrechnen und Terme (nach deren Prüfungen): Auftrag an die KI, Folgefehler, (E)-Zeilen.
\documentclass[11pt]{article}
\usepackage{../lp-druck}
\renewcommand{\lpname}{Bruchterme, Gleichungen und Formeln}
\renewcommand{\lpteilgebiet}{Gleichungen ·}
\renewcommand{\lpfuss}{Bewertungspaket}
\newcolumntype{P}{>{\centering\arraybackslash}p{10mm}}
\newenvironment{raster}{\par\smallskip\noindent\renewcommand{\arraystretch}{1.15}\linespread{1.35}\selectfont\tabularx{\linewidth}{@{}>{\raggedright\arraybackslash}p{38mm}P X@{}}\toprule
  \small\textsc{Teilschritt} & \small\textsc{P} & \small\textsc{Musterlösung}\\\midrule}{\bottomrule\endtabularx\par}
\newcommand{\fehler}[1]{\par\smallskip{\small\textcolor{lprot}{\bfseries Typische Fehler:} #1}}
\newcommand{\E}{\,{\footnotesize\bfseries\color{lpgruen}(E)}}
\begin{document}
\lpkopf{Bewertungspaket zum Gesamttest}{}

\begin{lpachtung}{Erst lösen, dann öffnen}
Dieses Paket enthält die Musterlösung. Wer es vor dem Lösen liest, misst am Ende nur, wie gut das Abschreiben gelingt.
\end{lpachtung}

\section*{1 · So gehst du vor}
\begin{enumerate}[leftmargin=*,itemsep=2pt]
\item \textbf{Gesamttest lösen} — vollständig, mit Rechenweg, ohne dieses Paket und ohne Taschenrechner.
\item \textbf{Selbst bewerten} mit Abschnitt~3: Zeile für Zeile vergleichen und die Punkte eintragen. Das geht ohne Hilfsmittel
  und ist der Normalfall.
\item \textbf{Oder bewerten lassen.} Ob du dafür eine KI nutzt und welche, entscheidest du selbst und in eigener Verantwortung —
  je nachdem, was dir aufgrund deines Alters und deiner persönlichen Situation erlaubt ist (Altersgrenzen und
  Nutzungsbedingungen des Anbieters, allenfalls Einverständnis deiner Eltern). Fotografiere oder scanne deine Lösung gut lesbar,
  \textbf{ohne Namen}. Handyfotos tragen oft unsichtbar Standort, Zeit und Gerät mit (Metadaten) — lade lieber einen Scan oder ein
  Bildschirmfoto des Fotos hoch. Lade nur deine Lösung, das Testblatt und dieses PDF hoch, keine persönlichen Daten, und füge den Auftrag aus
  Abschnitt~2 ein.
\item \textbf{Rückmeldung prüfen.} Stimmt, was die KI aus deiner Handschrift gelesen hat? Eine KI kann sich irren —
  im Zweifel gilt das Raster in Abschnitt~3.
\item \textbf{Gezielt wiederholen} nach Abschnitt~4.
\end{enumerate}

\section*{2 · Auftrag an die KI}
\begin{tcolorbox}[colback=lplinie!25,colframe=lplinie,boxrule=0.5pt,arc=1mm,fontupper=\small]
Du bist Korrektorin oder Korrektor für Mathematik auf der Sekundarstufe I in der Schweiz. Im Anhang findest du (1)~das Testblatt mit den Aufgaben, (2)~das Bewertungspaket zum Gesamttest «Bruchterme, Gleichungen und Formeln» mit Musterlösung und Punkteraster und (3)~Fotos meiner handschriftlichen Lösung.\par\medskip
Geh so vor:\par
1. Schreib zu jeder Aufgabe G1 bis G9 zuerst ab, was du in meiner Lösung liest — Rechenweg und Ergebnis. Was du nicht sicher lesen kannst, markierst du mit [unsicher]. Nicht raten. Achte besonders auf Minuszeichen, Hochzahlen und Klammern.\par
2. Vergib die Punkte genau nach dem Raster im Bewertungspaket (Abschnitt 3), ganze Punkte je Zeile. Ziehe einen Fehler nur einmal ab: Rechne ich mit einem falschen Zwischenergebnis richtig weiter, gibt es die Folgepunkte. Ausnahme: Zeilen mit \textbf{(E)} gibt es nur für das richtige Ergebnis, nie als Folgepunkt. Die Zeilen «Typische Fehler» sagen, welche Rasterzeile bei diesem Fehler 0 Punkte gibt — sie sind kein zusätzlicher Abzug.\par
3. Andere richtige Lösungswege zählen voll, zum Beispiel eine andere Reihenfolge der Umformungen, ein anderer gemeinsamer Nenner oder Zahlen erst eingesetzt und dann umgeformt (G8) — \textbf{ausser} die Aufgabe verlangt ein Verfahren ausdrücklich: G4 die Probe, G7(a) und G8 die umgestellte Formel. Gleichwertig: $\tfrac{x - 5}{3}$ und $\tfrac{1}{3}(x - 5)$; $L = \{4\}$ und «$x = 4$»; $L = \mathbb{Q}$ und «alle Zahlen». Ein Bruchterm zählt nur \textbf{vollständig gekürzt} und auf \textbf{einem Bruchstrich}.\par
4. Begründe jeden Abzug in einem Satz und nenne die Fehlerart (zum Beispiel «Summanden gekürzt», «Nenner addiert», «Minus vor dem Bruch», «Umkehrung verwechselt»). Schreib die Musterlösung nicht ab — zeig mir, wo mein Fehler liegt.\par
5. Gib zum Schluss eine Tabelle «Aufgabe | Punkte | Begründung», die Gesamtpunktzahl (von 25) und — nach der Selbsteinschätzung in Abschnitt 4 — welches Kapitel des Leitprogramms ich wiederholen soll. Keine Note.\par
6. Fehlt ein Foto oder ist eine Aufgabe unlesbar, sag es, statt sie zu bewerten.
\end{tcolorbox}

\needspace{16\baselineskip}
\section*{3 · Musterlösung und Punkteraster}
\textbf{Allgemein:} Ganze Punkte je Zeile. Ein Fehler wird nur einmal abgezogen (Folgefehler). Gleichwertige Wege zählen voll.
In der Spalte P: \textbf{(E)}~= Ergebnispunkt, Ergebnis richtig und vollständig gekürzt \emph{und} Weg erkennbar; ohne Zeichen = die
Zeile bewertet den Weg oder die Begründung. Der ganze Test ist ohne Taschenrechner gestellt.
«Typische Fehler» nennt, welche Zeile 0 Punkte gibt, und die Punkte, die bleiben.
\textbf{Teilrichtige Zeilen:} Verlangt eine Zeile mehrere Angaben und stimmt nur ein Teil, gibt die Zeile 0 Punkte,
ausser die Zeile nennt ausdrücklich eine Aufteilung.
\textbf{Zahlen:} Dezimalzahl oder Bruch, wenn gleich ($2.5 = \tfrac{5}{2}$). Einheiten in G7 und G8 sind nicht verlangt.

\begin{lpaufgabe}{G1}{Kürzen}{3 P}
\begin{raster}
Zerlegen & 1 & Zähler $(x - 5)^2$ (binomische Formel), Nenner $3(x - 5)$ (ausklammern)\\
Ergebnis & 1\E & $\dfrac{(x - 5)^2}{3(x - 5)} = \dfrac{x - 5}{3}$\\
Nicht erlaubt & 1 & $x = 5$ (Nenner $3x - 15 = 0$)\\
\end{raster}
\fehler{Summanden gekürzt (z.\,B. $\tfrac{x^2 + 25}{3x}$ oder $\tfrac{x}{3} - \ldots$): Zerlege- und Ergebniszeile 0; die nicht erlaubte Zahl gibt ihren Punkt $\to$ 1 von 3\,P.
$x = -5$ statt $x = 5$: Zeile «nicht erlaubt» 0 $\to$ 2 von 3\,P.
Richtig zerlegt, aber nicht gekürzt: Ergebniszeile 0 $\to$ 2 von 3\,P (mit richtiger Zahl).}
\end{lpaufgabe}

\begin{lpaufgabe}{G2}{Addieren}{3 P}
\begin{raster}
Hauptnenner & 1 & $x(x - 1)$ (jeder andere gemeinsame Nenner zählt, wenn er richtig verwendet wird)\\
Erweitern & 1 & $\dfrac{3(x - 1) + 2x}{x(x - 1)}$\\
Ergebnis & 1\E & $\dfrac{5x - 3}{x(x - 1)}$ (gleichwertig: $\dfrac{5x - 3}{x^2 - x}$)\\
\end{raster}
\fehler{Nenner addiert ($\tfrac{5}{2x - 1}$): alle Zeilen 0.
Zähler nicht erweitert ($\tfrac{5}{x(x - 1)}$): Erweitern- und Ergebniszeile 0 $\to$ 1 von 3\,P.
$3(x - 1) = 3x - 1$: Erweitern- und Ergebniszeile 0 $\to$ 1 von 3\,P.}
\end{lpaufgabe}

\begin{lpaufgabe}{G3}{Fehler finden}{3 P}
\begin{raster}
Probe & 1 & Mit $x = 1$: links $\dfrac{10}{4} = 2.5$, rechts $1 + 3 = 4$. Nicht gleich.\\
Fehler & 1 & Gekürzt wurde, als ob $x^2 + 9 = (x + 3)^2$ wäre, oder es wurden Summanden gekürzt; Kürzen geht nur mit gemeinsamen Faktoren.\\
Kürzbar? & 1 & Nein: $x^2 + 9$ lässt sich nicht in Faktoren zerlegen; es gibt keinen gemeinsamen Faktor mit $x + 3$.\\
\end{raster}
\fehler{Nur «falsch» ohne Probe und ohne Grund: 0 von 3\,P.
«Man darf nicht kürzen» ohne Hinweis auf Faktoren/Summanden: Fehlerzeile 0.
$x^2 + 9 = (x + 3)(x - 3)$ behauptet: Zeilen «Fehler» und «kürzbar?» 0 $\to$ höchstens 1 von 3\,P.}
\end{lpaufgabe}

\begin{lpaufgabe}{G4}{Klammern}{3 P}
\begin{raster}
Klammer & 1 & $7x - 3x + 6 = 2x + 14$: Das Minus vor der Klammer dreht beide Vorzeichen.\\
Lösung & 1\E & $4x + 6 = 2x + 14$, $2x = 8$, $x = 4$; $L = \{4\}$\\
Probe & 1 & $28 - 3 \cdot 2 = 22$ und $8 + 14 = 22$. Die Probe zählt auch als Folgepunkt, wenn sie mit dem eigenen Wert richtig gerechnet ist und den Fehler zeigt.\\
\end{raster}
\fehler{$-3(x - 2) = -3x - 6$: Klammer- und Lösungszeile 0 ($x = 10$); Probe richtig gerechnet $\to$ 1 von 3\,P.
Probe fehlt: Probezeile 0 $\to$ höchstens 2 von 3\,P.}
\end{lpaufgabe}

\begin{lpaufgabe}{G5}{Brüche}{3 P}
\begin{raster}
Bruchfrei & 1 & Mal $12$: $3(x + 3) - 2(x - 1) = 24$, beide Seiten multipliziert, Zähler in Klammern\\
Klammern & 1 & $3x + 9 - 2x + 2 = 24$\\
Lösung & 1\E & $x + 11 = 24$, $x = 13$; $L = \{13\}$\\
\end{raster}
\fehler{Rechte Seite nicht multipliziert ($3(x + 3) - 2(x - 1) = 2$, $x = -9$): erste Zeile und Lösungszeile 0 $\to$ 1 von 3\,P.
Minus vor dem Bruch nur beim ersten Glied ($- 2x - 2$, $x = 17$): Klammer- und Lösungszeile 0 $\to$ 1 von 3\,P.}
\end{lpaufgabe}

\begin{lpaufgabe}{G6}{Lösungsmenge}{2 P}
\begin{raster}
Umformen & 1 & $4x - 4 + 6 = 4x + 2$, also $4x + 2 = 4x + 2$; die $x$ fallen weg, es bleibt $2 = 2$.\\
Lösungsmenge & 1\E & Wahr für jedes $x$: $L = \mathbb{Q}$ (alle Zahlen).\\
\end{raster}
\fehler{«$L = \{\,\}$» oder «$x = 0$»: Lösungsmengenzeile 0 $\to$ 1 von 2\,P, wenn richtig umgeformt.}
\end{lpaufgabe}

\begin{lpaufgabe}{G7}{Treppe}{3 P}
\begin{raster}
(a) Umstellen & 1 & $2h = 63 - a$ (zuerst $-a$)\\
(a) Ergebnis & 1\E & $h = \dfrac{63 - a}{2}$\\
(b) & 1 & $h = \dfrac{63 - 25}{2} = 19$ cm. Auch als Folgepunkt mit der eigenen Formel, wenn richtig eingesetzt.\\
\end{raster}
\fehler{$h = \tfrac{63}{2} - a$ ($= 6.5$ bei $a = 25$): beide Zeilen (a) 0; (b) als Folgepunkt 1 $\to$ 1 von 3\,P.
$h = 63 - 2a$: (a) beide 0 $\to$ höchstens 1 von 3\,P.}
\end{lpaufgabe}

\begin{lpaufgabe}{G8}{Zins}{3 P}
\begin{raster}
Umstellen & 1 & Mal $100$: $100Z = K \cdot p \cdot t$\\
Formel & 1 & $p = \dfrac{100 Z}{K \cdot t}$ (durch das ganze Produkt $K \cdot t$)\\
Wert & 1\E & $p = \dfrac{100 \cdot 150}{2000 \cdot 3} = \dfrac{15000}{6000} = 2.5$ (Prozent)\\
\end{raster}
\fehler{$p = \tfrac{100Z}{K} \cdot t$ (nur durch $K$ geteilt): Formelzeile 0; Wert $22.5$ gibt 0 $\to$ 1 von 3\,P.
Zuerst Zahlen eingesetzt und dann richtig nach $p$ aufgelöst: gleichwertig, volle Punkte.}
\end{lpaufgabe}

\begin{lpaufgabe}{G9}{Fehler finden}{2 P}
\begin{raster}
Probe & 1 & Mit $G = 12$, $h = 5$ ist $V = 20$. Die falsche Formel gibt $\dfrac{20 \cdot 5}{3} = \dfrac{100}{3}$, nicht $12$.\\
Richtig & 1\E & Mal $3$, dann durch $h$: $G = \dfrac{3V}{h}$ (Probe: $\dfrac{60}{5} = 12$)\\
\end{raster}
\fehler{Nur «falsch» ohne Probe: Probezeile 0.
$G = \tfrac{V}{3h}$: Richtigzeile 0 $\to$ 1 von 2\,P, wenn die Probe stimmt.}
\end{lpaufgabe}

\needspace{12\baselineskip}
\section*{4 · Selbsteinschätzung}
\begin{tabularx}{\linewidth}{@{}l X@{}}
\toprule
\textbf{22 – 25 P} & Die geprüften Teile sitzen. Wo du Punkte verloren hast: das Kapitel dieser Aufgabe nochmals.\\
\textbf{17 – 21 P} & Den schwächsten Teil nochmals, gemessen am Anteil (Teil A 9, B 8, C 8 P): Tüfteln und Übungen der Kapitel dieses Teils.\\
\textbf{11 – 16 P} & Zurück zu den Kapiteln aller Aufgaben, in denen du Punkte verloren hast.\\
\textbf{0 – 10 P} & Zurück zu den Kapiteln jedes Teils, in dem du weniger als die Hälfte erreicht hast; sind es alle drei, zum Vortest und zu Kapitel 1.\\
\bottomrule
\end{tabularx}

\smallskip
\textbf{Aufgabe $\to$ Kapitel:} G1, G3 $\to$ 1; G2 $\to$ 2; G4 bis G6 $\to$ 3; G7 bis G9 $\to$ 4.
\end{document}
